% ch-introduction.tex — Introduction chapter

%#
%## CHAPTER: Introduction
%#
\chapter{Introduction}\label{cha:introduction}

This is the introductory chapter of the monograph.
It provides an overview of the main topics and motivates the central questions.

%#
%### SECTION: Background and Motivation
%#
\section{Background and Motivation}\label{sec:background-motivation}

The study of \ldots{} has a long history in mathematics.
We begin by recalling the fundamental definitions.

\begin{definition}\label{def:fundamental-concept}
  A \emph{widget} is a pair $(X, \tau)$ where $X$ is a set and
  $\tau \subseteq \mathcal{P}(X)$ satisfies the following axioms:
  \begin{enumerate}
  \item
    $\emptyset, X \in \tau$;
  \item
    $\tau$ is closed under arbitrary unions;
  \item
    $\tau$ is closed under finite intersections.
  \end{enumerate}
\end{definition}

%#
%### SECTION: Overview of Results
%#
\section{Overview of Results}\label{sec:overview-results}

The main result of this monograph is the following.

\begin{theorem}\label{thm:main-result}
  Every compact widget is sequentially compact.
\end{theorem}

The proof of \cref{thm:main-result} occupies \cref{cha:main-results}.

%#
%### SECTION: Organization
%#
\section{Organization}\label{sec:organization}

The monograph is organized as follows.
\Cref{cha:background} reviews the necessary prerequisites.
\Cref{cha:main-results} contains the proof of the main theorem.
\Cref{cha:conclusion} discusses open problems and future directions.

%%% Local Variables:
%%% mode: LaTeX
%%% TeX-master: "monograph-template"
%%% End:
